\section{Methods}
\label{sec:methods}

\subsection{Data}
\label{sec:methods-data}

\paragraph{Corpus.} Tahoe-100M~\cite{tahoe100m}: drug-treated single-cell transcriptomes across
cancer cell lines, with plate-matched DMSO controls. Metadata supplies cell line
(Cellosaurus ID), drug, dose (\texttt{drugname\_drugconc}), plate, and a fine-grained
mechanism-of-action annotation (\texttt{moa-fine}).

\paragraph{Gene panel.} All analysis is restricted to the $946$ genes in the intersection of the
L1000 landmark set~\cite{l1000} with Tahoe's measured genes. Expression is library-normalised to
counts-per-10k and $\log1p$-transformed.

\paragraph{A representational pitfall worth recording.} Tahoe's per-cell \texttt{genes} field holds
\emph{vocabulary token identifiers}, not positional indices into the gene metadata table. Mapping
them positionally scrambles gene identity silently --- values exceed the metadata row count, so the
error surfaces only as an out-of-bounds index at the tail. All expression construction here maps
\texttt{gene\_metadata['token\_id']} $\rightarrow$ gene symbol $\rightarrow$ panel index.
Validation: a PCA cell-line probe on the resulting panel reaches $0.864$ against a chance rate of
$0.02$.

\paragraph{Unit of analysis.} A \emph{condition} is a (drug, cell line, plate, dose) tuple. Plates
are numbered within cell line rather than globally, so the batch unit is the pair (cell line, plate)
and never the plate label alone --- a distinction that matters for every within-plate comparison
below.

\paragraph{Scale.} Table~\ref{tab:scale} states the quantities every $n$ in this thesis should be
located against. Two populations are used and they are not interchangeable. The broad
characterisation of the task (Section~\ref{sec:res-blind}) streams a wide sample of the atlas; the
residual arm (Sections~\ref{sec:res-encoding} onward) works from a smaller cached subset built once
and reused, so that every arm sees identical cells.

\begin{table}[h]
\centering
\begin{tabular}{lr}
\toprule
\multicolumn{2}{l}{\textbf{Task characterisation} (streamed sample)} \\
\midrule
conditions characterised & 6,628 \\
\quad of which statistically active against DMSO & 3,647 / 5,014 tested (72.7\%) \\
\quad identifiable in principle (ceiling $\geq 0.8$) & 3,064 (46.2\%) \\
\quad with a plate-mate closer than their own replicate & 78.7\% \\
median cells per condition & 44 \\
median \#DEGs per condition (BH $q<0.05$) & 0 \\
median SNR (effect $\div$ replicate noise) & 0.75 \\
drugs with a real dose series & 218 / 351 (62.1\%) \\
drugs seen in $\geq 3$ cell lines & 227 \\
\midrule
\multicolumn{2}{l}{\textbf{Residual-frame cache}} \\
\midrule
cells & 620,564 (600,000 treated, 20,564 control) \\
(cell line, plate) groups & 198 \\
drugs / cell lines & 106 / 50 \\
conditions with $\geq40$ treated and $\geq20$ control cells & 6,617 \\
\quad surviving the reliability filter ($\cos(r_A,r_B) > 0.2$) & 4,091 (61.8\%) \\
training examples (60 control cells per condition) & 245,444 \\
\midrule
\multicolumn{2}{l}{\textbf{Evaluation sets}} \\
\midrule
tier 1 (seen conditions) / tier 2 (unseen drugs) / tier 3 (unseen combinations) & 4,325 / 2,188 / 99 \\
conditions scored per tier in the generation evaluations & 400 \\
tier-2 drugs in the expression-frame $\NIR$ benchmark (within plate) & 606 \\
conditions scored in the residual-frame evaluation & 570 (41 cell lines, 90 drugs) \\
\bottomrule
\end{tabular}
\caption{The populations behind every $n$ quoted in this thesis. Sources:
\texttt{drug\_biology\_atlas.json}, \texttt{eval\_endcell\_cpu.json},
\texttt{nir\_sameplate.json}, \texttt{re\_holdout2\_stratified.json},
\texttt{target\_divergence.json}. Note that the 72.7\% active figure is over the 5,014 conditions
with enough cells for the permutation test, not over all 6,628.}
\label{tab:scale}
\end{table}

\paragraph{Splits.} Three evaluation tiers, defined by what the training set contained.
\emph{Tier 1} holds conditions whose (drug, cell line) pairing was seen in training, so it measures
fit rather than generalisation. \emph{Tier 2} holds conditions whose \emph{drug} never appeared,
which is Leave-Drugs-Out in the vocabulary of Section~\ref{sec:lit-dreval}. \emph{Tier 3} holds
pairings whose drug and cell line were each seen but never together --- Leave-Pairs-Out. Tier 3 is
small ($99$ conditions) and is reported only where it is adequately powered; the residual arm builds
its own, larger, (drug, cell line)-keyed holdout for the same purpose
(Section~\ref{sec:methods-holdout}).

\subsection{Model and representation}

\paragraph{Cell sentences.} A cell is written as its expressed panel genes in descending order of
expression, terminated by \ENDCELL. The prompt supplies cell line, drug, dose, mechanism, and the
untreated control cell's own sentence; the response is the treated cell's sentence.

\paragraph{Backbone.} \texttt{C2S-Scale-Pythia-1b}~\cite{c2sscale}, cold-started for every arm from
the same base checkpoint with identical hyperparameters (1 epoch, learning rate $10^{-5}$,
gradient accumulation 16, maximum length 8192, bf16), so that in each experiment the training
target or objective is the only variable.

\paragraph{Sentinels.} \ENDCELL{} and, for the residual arm, \DOWN{} are registered as atomic
special tokens. This matters: unregistered, \DOWN{} splits into sub-words and the up/down boundary
ceases to be a clean symbol.

\subsection{Metrics}
\label{sec:methods-metrics}

\paragraph{The standard metric.} The field's default is a correlation between predicted and true
\emph{change from control}, restricted to the genes that actually move. Writing $x$ for the control
profile, $y$ for the treated truth and $\hat{y}$ for the prediction, let
\begin{equation}
  S_K \;=\; \operatorname*{arg\,top-}K_{\,g} \; \lvert y_g - x_g \rvert
\end{equation}
be the $K$ genes with the largest true change. Then
\begin{equation}
  \DEdr \;=\; \operatorname{corr}\!\big( \left(\hat{y} - x\right)\big|_{S_K}, \;
                                          \left(y - x\right)\big|_{S_K} \big).
\end{equation}
We report $K=50$ with Pearson correlation throughout; $K \in \{20, 100, 200\}$ and the Spearman
variant were computed and move the numbers by less than $0.02$ except at $K=200$, where the
top-$K$ set starts to include genes that do not move.

Three details are load-bearing rather than incidental, because Section~\ref{sec:res-metrics} turns
on them.

First, $S_K$ is selected using the \emph{truth}, so every predictor is scored on the genes that
condition actually moved. Second, and consequently, $\DEdr$ rewards agreement with $y - x$ and a
prediction that simply reverts toward the population centre correlates with $y - x$ by construction
--- which is the artifact the results chapter exhibits. Third, \textbf{partial-}$\DEdr$ is the same
quantity with the control profile regressed out of both arguments before correlating, which removes
that route and is therefore the honest version.

\paragraph{Absent genes.} A cell sentence lists only expressed genes, so a predicted profile leaves
some panel genes unranked. Two conventions were used throughout and reported side by side:
\texttt{worst} assigns every absent gene the last rank $P$, and \texttt{francesca} assigns them a
fixed middle rank $P/2$. The choice is immaterial --- spike-in $\DEdr$ reads $0.952$ against $0.954$
under the two, and no baseline or ceiling in this thesis moves by more than $0.01$ --- so all figures
below are quoted under \texttt{worst} and the convention is not mentioned again.

\paragraph{$\paneltau$.} Kendall's $\tau$ between predicted and true rank over the full $946$-gene
panel. Unlike $\DEdr$ it selects no gene subset and is invariant to the absent-gene convention by
construction.

\paragraph{Normalised Inverse Rank.} Introduced by Wu et al.\ in
PerturBench~\cite{perturbench} (there called the ``transposed-rank'' metric) and identified by
Miller et al.~\cite{miller} as one of the few consistently calibrated metrics across 14 Perturb-seq
datasets. For a prediction $\hat{y}$ of condition $i$, with truths $\{y_j\}$ over the comparison
set,
\begin{equation}
  \NIR(\hat{y}, i) \;=\; \frac{1}{|J|}\sum_{j \in J}
    \Big[\mathbf{1}\{\sigma(\hat{y}, y_i) > \sigma(\hat{y}, y_j)\}
       + \tfrac{1}{2}\,\mathbf{1}\{\sigma(\hat{y}, y_i) = \sigma(\hat{y}, y_j)\}\Big],
\end{equation}
the fraction of \emph{other} conditions whose truth the prediction resembles less than its own,
with ties counted at one half. Chance is $0.50$. Because the generic response is present in every
condition's truth in equal measure, it cancels in the comparison: $\NIR$ is a
\emph{discrimination} metric, and only the drug-specific difference decides it.

The tie term is not cosmetic. In residual space (Section~\ref{sec:methods-residual}) a predictor of
``the average drug response'' is exactly the zero vector, every similarity against it is zero, and a
strict inequality scores it $0.000$ rather than the $0.500$ it deserves.

\paragraph{Which $\NIR$ variant, and why it must be stated.} Our instruments compute $\NIR$ under
two similarity functions and store both: \texttt{nir\_expr}, Euclidean distance in expression space,
and \texttt{nir\_rank}, the same construction over rank vectors. \textbf{Every $\NIR$ figure in this
thesis is \texttt{nir\_expr}}, and the choice is load-bearing rather than incidental: under
\texttt{nir\_rank} the tier-2 aggregate places \emph{every} arm below chance, including the noise
ceiling itself ($0.380$, against the model's $0.313$ and a drug-agnostic linear map's $0.285$). A
metric on which a real replicate of a condition fails to be retrieved above chance is not measuring
what its name suggests, and it is not the metric the calibration analysis of
Section~\ref{sec:res-metrics} endorsed --- the $\DRF$ result is computed on the expression-space
variant. We report the calibrated one and name it here so the stored artifacts cannot be misread.

% fig-nir.tex -- NIR as a discrimination score. Goes with the metric definition.
% No data dependency.
\begin{figure}[htbp]
\centering
\begin{tikzpicture}[
  x=1cm, y=1cm, font=\small,
  lbl/.style={font=\scriptsize, text=black!70},
  strong/.style={font=\scriptsize, text=black!90},
  own/.style={circle, fill=black!85, inner sep=2pt},
  other/.style={circle, draw=black!55, fill=white, inner sep=1.9pt},
  ar/.style={-{Stealth[length=4.5pt]}, black!65},
]

\node[lbl, anchor=south] at (2.05,0.44) {other drugs in the same comparison set};
\node[strong, anchor=south] at (6.30,0.44) {own truth};

\fill[black!10, rounded corners=1pt] (0.55,0.00) rectangle (6.30,0.34);
\foreach \x in {0.95, 1.75, 2.60, 3.55, 4.45, 5.45} {\node[other] at (\x,0.17) {};}
\node[own] at (6.30,0.17) {};
\node[other] at (7.15,0.17) {};
\node[other] at (7.85,0.17) {};
\node[anchor=west, font=\small] at (8.35,0.17) {$\NIR = \tfrac{6}{8}$};

\draw[ar] (0.45,-0.26) -- (8.15,-0.26);
\node[lbl, anchor=north] at (4.30,-0.38)
  {similarity of the prediction to each condition's truth $\longrightarrow$};

\draw[decorate, decoration={brace, mirror, amplitude=3.5pt}, black!55]
  (0.55,-0.80) -- (6.30,-0.80);
\node[strong, anchor=north] at (3.40,-0.96) {the shaded fraction \emph{is} the score};

\node[lbl, anchor=west] at (0.45,-1.50) {own truth ranked first $\Rightarrow \NIR \to 1$};
\node[lbl, anchor=west] at (0.45,-1.84) {own truth ranked at random $\Rightarrow \NIR = 0.5$};
\node[strong, anchor=west] at (0.45,-2.18)
  {every similarity \emph{identical} $\Rightarrow \NIR = 0.5$ by the half-tie rule, not $0$};

\end{tikzpicture}
\caption[NIR is a discrimination score]{\textbf{$\NIR$ asks whether a prediction resembles its own
condition's truth more than it resembles the truths of the other drugs in the same comparison set.}
The generic response --- the stress and cell-cycle program every compound triggers --- is present in
every truth in equal measure, so it cancels in the comparison and only the drug-specific difference
decides. Two consequences follow, and both matter downstream. Chance at $0.50$ is a statement about
the \emph{comparison set}, not about the model: conditions are grouped by (cell line, plate) in the
expression frame and by cell line in the residual frame, and the two are not interchangeable. And a
predictor whose similarities are all identical --- the zero vector, which is what ``predict the
average drug response'' becomes in residual space --- must score $0.50$ rather than $0$, which is why
ties are counted at one half rather than dropped.}
\label{fig:nir}
\end{figure}


\paragraph{Metric calibration (\DRF).} We port the framework of Miller et al.~\cite{miller}, in
which each \emph{metric} is graded rather than each model, by asking how well it separates a
positive from a negative control:
\begin{equation}
  \DRF(m) \;=\; \frac{m(\textrm{pos}) - m(\textrm{neg})}{m(\textrm{perfect}) - m(\textrm{neg})},
\end{equation}
with $\textrm{pos}$ an interpolated-duplicate noise ceiling --- their positive control, which
interpolates a technical duplicate toward the mean baseline with per-gene weights given by
differential-expression significance --- and $\textrm{neg}$ a drug-agnostic leave-one-out mean. A
negative \DRF{} means the metric scores the uninformative baseline \emph{above} the noise ceiling:
it is inverted, and rewards the generic gene ordering rather than the drug.

Adopting this framework is deliberate. Miller et al.\ use it to argue that reports of deep models
failing to beat uninformative baselines reflect miscalibrated evaluation rather than model failure
(Section~\ref{sec:lit-dispute}). Evaluating our model on the metric \emph{their} analysis endorses
means the result reported in Section~\ref{sec:res-blind} cannot be dismissed on the grounds their
own paper raises.

\subsection{Controls}
\label{sec:methods-controls}

The negative results in this thesis are only as strong as the controls that make them
interpretable. Each of the following was added in response to a specific objection.

\paragraph{Plate confounding.} Drug and plate are partially confounded by Tahoe's design: each drug
is assayed on its own plate(s), with roughly 4--12 drugs per plate and plate-matched controls.
Discrimination grouped by cell line alone therefore lets the \emph{batch} signature stand in for
drug identity. Measured directly, by scoring the same tier-2 conditions with and without the plate
held constant:

\begin{table}[h]
\centering
\begin{tabular}{lrr}
\toprule
Predictor (drug-agnostic) & cross-plate & within plate \\
\midrule
mean baseline & 0.399 & \textbf{0.161} \\
control-copy & 0.551 & 0.510 \\
\midrule
noise ceiling & 0.625 & 0.575 \\
fraction of drugs at chance & 43\% & 49\% \\
\bottomrule
\end{tabular}
\caption{Expression-frame $\NIR$. A drug-agnostic mean baseline more than doubles its score when the
plate is free to vary, which is what plate leakage looks like. The ceiling falls once the plate is
held constant because the batch signature was inflating it too.}
\label{tab:plateleak}
\end{table}

The mean baseline is the clearest demonstration: it contains no drug information whatsoever, and it
scores $0.399$ when the comparison set may span plates against $0.161$ when it may not. Every
discrimination instrument therefore gained a \texttt{--same\_plate\_only} mode and all discrimination
results in this thesis were re-run within plate. Conclusions are unchanged; some magnitudes shrink.

A separate and stronger observation, which belongs to the drug-blindness argument rather than to this
one, is that \emph{within} plate and on the identifiable stratum a control-copy baseline reaches
$0.766$ against the model's $0.768$ (Section~\ref{sec:res-blind}). That is not a statement about
plate leakage --- it is a statement about the model.

\paragraph{The scramble ablation.} The central manipulation. The prompt is left byte-identical
except that the drug name and mechanism are replaced by another drug's, while the control cell and
the scoring truth are unchanged. Because both arms share the same control cell, control- and
batch-derived leakage cancels in the \emph{difference} $\gap = \NIR_{\textrm{model}} -
\NIR_{\textrm{scramble}}$. This is the only manipulation that isolates drug use, and every drug-use
claim in this thesis is anchored to it rather than to a comparison against a drug-agnostic
predictor.

\paragraph{Two further drug-isolating instruments.} The drug-blindness claim rests on four
instruments rather than one, because they fail in different ways and agreement across them is worth
more than any single number. Two are defined above ($\NIR$ and the scramble); the other two are:

\emph{Forced-choice grading.} A prediction is scored against its own condition's truth and against
another drug's truth from the same comparison set, and the instrument records how often the own
truth wins. Chance is $0.50$ by construction. The ceiling is the same forced choice with a real
held-out replicate substituted for the prediction, which bounds what any predictor could achieve at
this noise level; it lands between $0.67$ and $0.83$ depending on cells per condition.

\emph{Output invariance.} This asks whether the drug token changes the generated text at all,
without reference to any truth. For a fixed control cell, the model generates twice from the
\emph{same} prompt and once from each of several prompts differing only in the drug, and the
instrument compares
\begin{equation}
  \textrm{gap} \;=\; \underbrace{\overline{\textrm{sim}}(\textrm{same prompt, resampled})}_{\textrm{sampling noise alone}}
  \;-\; \underbrace{\overline{\textrm{sim}}(\textrm{prompts differing in the drug})}_{\textrm{noise plus any drug effect}},
\end{equation}
with similarity measured as top-$100$ Kendall $\tau$ and as Jaccard overlap. A positive gap means two
generations for the same drug resemble each other more than generations for different drugs do,
i.e.\ the drug token moves the output. A gap of zero means swapping the drug perturbs the output no
more than resampling at the same temperature does. Reported over $8$ contexts $\times$ $12$ drugs at
temperature $0.8$. This instrument is the weakest of the four --- it can only detect whether the
output moves, not whether it moves \emph{correctly} --- but it is also the one least able to be
confounded, since it needs no ground truth at all.

\paragraph{Swap similarity --- stratified scramble.} A scramble that swaps in a drug with a
\emph{similar} true response is a weak test: an unchanged output would then be correct, not blind.
This is a real risk here, because same-mechanism drugs are barely more similar than random pairs
(within-MoA distance $2.322$ versus between-MoA $2.377$, ratio $0.977$), so ``different mechanism''
does not imply ``different response''. Two remedies are used:
\begin{enumerate}
  \item a \emph{swap-distance sweep} that turns the single test into a curve of $\gap$ against the
        response distance $\lVert y_A - y_B \rVert$ between the real and swapped-in drug; and
  \item a \emph{stratified scramble} with three defined partners per condition chosen by residual
        cosine within the cell line --- \texttt{near} (most similar), \texttt{orth} ($\cos\approx0$)
        and \texttt{opposite} (most anti-correlated). Genuine drug use must produce a gap that
        \emph{grows} across the strata; a flat profile is a red flag.
\end{enumerate}

% fig-scramble.tex -- the scramble ablation and why its difference is leak-immune.
% No data dependency.
% Layout note: the prompt boxes are ~2.2cm tall once the shaded spans are set, which is far more
% than a four-line block looks like it should need. Every band below them is placed against that
% measured height, not against an estimate.
\begin{figure}[htbp]
\centering
\begin{tikzpicture}[
  x=1cm, y=1cm, font=\small,
  lbl/.style={font=\scriptsize, text=black!70},
  strong/.style={font=\scriptsize, text=black!90},
  promptbox/.style={draw=black!35, rounded corners=2pt, inner sep=5pt, fill=black!2},
  ar/.style={-{Stealth[length=4.5pt]}, black!65},
]
\newcommand{\spanhl}[1]{\colorbox{black!14}{#1}}

\node[lbl, anchor=west] at (0.15,0.00) {real prompt};
\node[promptbox, anchor=north west] (p1) at (0.15,-0.26) {%
  \renewcommand{\arraystretch}{0.9}%
  \begin{tabular}{@{}l@{~~}l@{}}
  \ttfamily\scriptsize Cell line:    & \ttfamily\scriptsize MCF7 \\
  \ttfamily\scriptsize Drug:         & \ttfamily\scriptsize \spanhl{Lapatinib} \\
  \ttfamily\scriptsize Mechanism:    & \ttfamily\scriptsize \spanhl{EGFR inhibitor} \\
  \ttfamily\scriptsize Control cell: & \ttfamily\scriptsize ACTB GAPDH \ldots \\
  \end{tabular}};

\node[lbl, anchor=west] at (0.15,-2.86) {scrambled prompt --- only the two shaded spans differ};
\node[promptbox, anchor=north west] (p2) at (0.15,-3.12) {%
  \renewcommand{\arraystretch}{0.9}%
  \begin{tabular}{@{}l@{~~}l@{}}
  \ttfamily\scriptsize Cell line:    & \ttfamily\scriptsize MCF7 \\
  \ttfamily\scriptsize Drug:         & \ttfamily\scriptsize \spanhl{Vorinostat} \\
  \ttfamily\scriptsize Mechanism:    & \ttfamily\scriptsize \spanhl{HDAC inhibitor} \\
  \ttfamily\scriptsize Control cell: & \ttfamily\scriptsize ACTB GAPDH \ldots \\
  \end{tabular}};

\node[draw=black!35, rounded corners=2pt, inner sep=6pt, align=center, fill=black!3]
  (truth) at (9.30,-2.55) {truth for\\\textbf{Lapatinib}\\{\scriptsize drawn once}};

\draw[ar] (p1.east) .. controls (6.6,-1.35) and (7.6,-1.85) .. (truth.north west);
\draw[ar] (p2.east) .. controls (6.6,-4.20) and (7.6,-3.30) .. (truth.south west);

\node[anchor=north, font=\small] at (9.30,-3.70)
  {$\gap = \NIR_{\text{real}} - \NIR_{\text{scrambled}}$};

\draw[ar, dashed, black!45] (1.20,-6.10) -- (8.10,-6.10);
\node[strong, anchor=south, align=center] at (4.65,-6.02)
  {the shared control cell enters \emph{both} arms, so control- and batch-derived};
\node[strong, anchor=north, align=center] at (4.65,-6.20)
  {leakage cancels in the difference};

\end{tikzpicture}
\caption[The scramble ablation]{\textbf{The scramble replaces the drug name and its mechanism and
nothing else.} Both arms are conditioned on the same control cell and scored against the same truth,
so any leakage through the control cell or the plate enters both arms in equal measure and cancels in
the difference $\gap$. This is why $\gap$, rather than a comparison against a drug-agnostic
predictor, is the quantity every drug-use claim in this thesis is anchored to: the drug-agnostic
comparison is not leak-immune, and a drug-agnostic mean baseline scores $\NIR = 0.399$ when the
comparison set may span plates against $0.161$ when it may not, precisely because it inherits that
leakage. A refinement matters here and is
developed in Section~\ref{sec:methods-controls}: if the swapped-in drug happens to have a
\emph{similar} true response then an unchanged output is correct rather than blind, so the swap
partner is chosen from three defined similarity strata and the evidence is the \emph{gradient} across
them rather than any single value.}
\label{fig:scramble}
\end{figure}


\paragraph{Field decomposition of the scramble.} The scramble above replaces the drug name
\emph{and} its mechanism together, so a null tells us the instruction was not read without saying
which part of it was ignored. Decomposing it costs one extra pair of generation arms and answers a
question the whole unseen-drug argument turns on. Each arm swaps exactly one span of the instruction
line to the opposite-signature partner and leaves every other byte identical: the drug name with the
mechanism kept, the mechanism with the name kept, and --- for continuity with earlier results --- both.

One guard is essential rather than tidy. If the swap partner happens to share the original's
mechanism, then swapping only the mechanism produces a prompt \emph{identical} to the model's own,
and the arm scores a gap of exactly zero \emph{by construction}. Read naively that is
indistinguishable from ``the model ignores the mechanism'', which is the very claim under test. Such
conditions are therefore dropped rather than scored, which is why the mechanism arm carries fewer
conditions than the others.

\paragraph{Channel gate.} For an unseen drug, prediction is only possible through a property the drug
\emph{shares} with drugs seen in training, so the reachable channels can be enumerated and each one
measured in the same frame. For a channel $X$, a drug's residual is predicted as the mean residual of
the \emph{other} drugs in the same cell line that share property $X$ with it --- the construction of
the mechanism baseline, generalised. Because it uses only drugs that were in training, it is a fair
contest in every split, unlike a same-drug lookup, which on an unseen drug is an oracle.

Two disciplines make the result readable. Every channel is paired with a \textbf{count-matched random
null}: the same \emph{number} of partner drugs, drawn at random from the same cell line and averaged
identically. Without it, a channel scoring above chance cannot be distinguished from the fact that
averaging $k$ residuals denoises, which is true for any $k$ and says nothing about the channel. And
\textbf{coverage is reported before any verdict}, because a channel annotated on few drugs cannot be
\emph{closed} by a null result --- it was never tested.

The bound this establishes is on \emph{similarity transfer} through each channel, not on any model.
A model could in principle combine channels or learn a non-smooth structure-to-response map that a
nearest-neighbour average cannot express. A leave-one-drug-out regression over all channels jointly
addresses the first. The second is bounded by sample size rather than by argument: at a few hundred
drugs, learning a non-smooth map is not feasible, so at this $n$ the lookup is effectively the bound.

\paragraph{Is the interaction learnable?} The transfer coefficient says how \emph{large} the
drug$\times$cell-line interaction is; it says nothing about whether anything could learn it. If a cell
line modulated drug responses in a consistent direction, that direction would be estimable from the
line's own cells and a conditional model would have a concrete target. Writing
$\kappa(d,c) = \resid(d,c) - \hat{\beta}(d)$, the test compares $\cos(\kappa(d,c), \kappa(d',c))$ for
different drugs in the \emph{same} cell line against the same quantity across \emph{different} cell
lines.

$\hat{\beta}$ must be estimated \textbf{leave-one-cell-line-out}. Estimated over all lines it would
contain $\resid(d,c)$ itself, and subtracting a mean containing its own term induces a $-1/(m-1)$
correlation that biases the signal arm downward \emph{by construction} --- which would make a
structured interaction look idiosyncratic, the exact error the test exists to avoid. Both arms are
disattenuated by $\kappa$'s own split-half reliability, since $\kappa$ is a difference of two noisy
quantities and is therefore noisier than the residual it comes from.

\paragraph{Uncertainty.} All confidence intervals are clustered bootstraps resampling \emph{cell
lines}, not conditions. Conditions within a cell line are pseudoreplicates in Hurlbert's
sense~\cite{hurlbert} --- multiple measurements of one unit treated as independent --- and an
observation-level interval would be badly anti-conservative. This is one of the six pathologies
DrEval~\cite{dreval} identifies in this field, and the structure of our data is the same as the one
they describe: hundreds of thousands of cells resolving to a few hundred conditions over $50$ cell
lines. For the transfer coefficient
(Section~\ref{sec:methods-vardecomp}) the cluster is the \emph{drug}, since the estimand is a
per-drug property.

\paragraph{Generation validity.} Every generation-based evaluation reports the \ENDCELL{}/\DOWN{}
completion rate, the fraction of emitted symbols that are valid panel genes, and the duplicate rate,
before any score is quoted. This is not hygiene: a $600$-token generation cap silently truncated
$26\%$ of generations in an early run and \emph{halved} the measured effect.

\subsection{Mechanistic probing}
\label{sec:methods-probe}

To distinguish ``cannot represent the drug'' from ``represents but does not use it'':

\paragraph{Decodability.} Per-layer linear probes on the residual stream at the last prompt token,
predicting drug identity.

\paragraph{Confound purification.} The raw drug subspace is partly a batch subspace --- probing
dose and plate from \emph{inside} it recovers dose at $0.90$--$0.95$. A context subspace is built
from the union of cell line, plate and dose, and the drug axis is orthogonalised against it. Every
causal claim uses the purified subspace.

\paragraph{Removal gate.} Before any ``ablating the drug does nothing'' result is trusted, the
subspace must be shown to contain the drug: the fraction of above-chance probe accuracy removed by
projecting it out must exceed $0.8$. Without this gate the null is unfalsifiable.

\paragraph{Causal test.} Log-probabilities of the teacher-forced true response are read clean, then
re-read with a forward hook projecting the subspace out at every position; the statistic is the mean
KL divergence over response tokens. Measuring in logit space rather than by generating removes the
sampling-noise floor that made an earlier steering experiment inconclusive.

\paragraph{The identity test.} Ablation shows a subspace carries functional variance; it cannot show
that the subspace's \emph{drug-specific content} is read. Only a swap can. For each cell that truly
received drug A, drug A's code is overwritten with drug B's \emph{within the purified subspace}, and
the resulting KL is compared against injecting a random vector of \emph{identical norm}. Drug means
are estimated on a held-out half of the cells and the swap performed on the disjoint half, so
mean-estimation noise cannot itself manufacture a null.

\subsection{Target constructions}

\paragraph{Consensus.} Cells grouped by (drug, cell line, dose) and averaged into one pseudobulk
sentence, with every original prompt retained so only the response changes.

\paragraph{Optimal transport.} For each (cell line, plate) group, a coupling $\pi$ between control
and treated cells minimising $\sum_{ij}\pi_{ij}\lVert x_i - y_j\rVert^2$ under uniform marginals,
solved by Sinkhorn iteration~\cite{cuturi} on $K = \exp(-C/\varepsilon)$. Couplings are computed in
the released Tahoe scVI latent~\cite{scvi} ($n_{\textrm{latent}} = 10$), joined per cell by barcode
rather than re-encoded --- a re-encode produced a degraded latent (cell-line probe $0.11$ versus
$0.975$ for the shipped latent). Targets are built \emph{decoder-free}, as convex combinations of
real treated cells in expression space, so they always lie on the data manifold.

The regularisation parameter defines a ladder whose endpoints are experiments we had already run:
$\varepsilon\rightarrow\infty$ gives a uniform coupling and hence the treated mean (consensus);
$\varepsilon\rightarrow 0$ gives near-hard assignment, approximately a single state-matched treated
cell. This is what makes the sweep conclusive rather than anecdotal.

\subsection{Residual targets}
\label{sec:methods-residual}

For each condition (drug, cell line, plate, dose) with at least $40$ treated and $20$ plate-matched
control cells:
\begin{equation}
  \resid(d,c) \;=\; \underbrace{\big(\bar{y}_{\textrm{treated}} - \bar{y}_{\textrm{control}}\big)}_{\shift(d,c)}
  \;-\; \underbrace{\frac{1}{|D_c|}\sum_{d' \in D_c} \shift(d',c)}_{\textrm{generic}(c)} .
\end{equation}
The control subtraction removes the cell's own state; the mean-over-drugs subtraction removes the
\emph{generic drug program} --- the stress and cell-cycle response common to all compounds, which
Section~\ref{sec:res-metrics} shows accounts for essentially all of the model's apparent skill. What
remains is what makes \emph{this} drug different.

\paragraph{Scope, and a justification that did not survive.} The generic is estimated per
\emph{cell line} rather than per plate. The original argument was reproducibility --- $62\%$ of
conditions reproduce across a half-split at cell-line scope against $19\%$ at plate scope --- on the
reasoning that a plate holds too few drugs to estimate a mean, so subtracting a noisy one injects
noise.

That comparison does not hold up, and Section~\ref{sec:res-transfer} retracts it. Both figures were
computed with a \emph{shared} control pseudobulk subtracted from both halves, which makes the halves
correlate through their common control error --- and it inflates cell-line reproducibility far more
than plate reproducibility, because the cell-line generic is itself less noisy. Measured honestly,
with disjoint control halves, the two are comparable: $20\%$ at plate scope against $16\%$ at
cell-line scope. The reproducibility argument therefore favours neither, and three other diagnostics
favour plate scope outright.

The targets used throughout this thesis were nonetheless built at cell-line scope, because that
decision was made before the diagnostics existed. The consequence is stated rather than hidden: the
residual targets retain plate-level batch structure. This does not touch $\gap$, which is a paired
contrast in which any component shared by both arms cancels, but it does mean a plate-scoped rebuild
is the obvious next construction.

\paragraph{Reliability filter.} A condition is kept only if its residual reproduces across a
half-split, $\cos(\resid_A, \resid_B) > 0.2$. This retains $4{,}091$ of $6{,}617$ conditions
($62\%$); the remainder are noise.

\paragraph{Encoding.} \texttt{<up genes> \DOWN{} <down genes> \ENDCELL}, $100$ genes per block,
ordered by signed residual magnitude. Sign is biology, so it is given a symbol rather than being
folded into a magnitude ranking.

% fig-frame.tex -- the residual frame and the encoding fork.
% No data dependency. Annotations from RESULTS/target_divergence.json, cellline scope:
%   full 34.59 / shift 111.42 / residual 118.25 ; ceilings 0.7420 / 0.7420 / 0.7432.
% Layout notes:
%   * two subtraction ROWS, not an L-shape (the L put the generic term on a band that collided
%     with its own caption).
%   * the caption for the generic sits BELOW both rows, not between them -- between them it ran
%     under the divider and into the right-hand column.
%   * row labels are centred on profiles spaced 1.65cm apart; at 1.45cm "generic(c)" and
%     "residual r(d,c)" touch.
\begin{figure}[htbp]
\centering
\begin{tikzpicture}[
  x=1cm, y=1cm, font=\small,
  bar/.style={fill=black!55, draw=none},
  faintbar/.style={fill=black!22, draw=none},
  gene/.style={draw=black!35, rounded corners=1pt, inner sep=1.8pt, font=\ttfamily\scriptsize},
  shared/.style={gene, fill=black!5, text=black!45},
  differs/.style={gene, fill=black!16, text=black, draw=black!70, thick},
  lbl/.style={font=\scriptsize, text=black!70},
  strong/.style={font=\scriptsize, text=black!90},
  op/.style={font=\normalsize, text=black!75},
]
\newcommand{\prof}[4]{%
  \foreach \h [count=\i from 0] in {#3} {
    \fill[#4] ($(#1,#2)+(\i*3.0pt,0)$) rectangle ++(2.3pt,\h pt);
  }%
}

% ================================================ LEFT: two subtractions
\node[lbl]    at (0.52,3.34) {treated};
\node[lbl]    at (2.17,3.34) {control};
\node[strong] at (3.90,3.34) {shift $\shift(d,c)$};
\prof{0.10}{2.62}{17,12,9,14,7,10,5,8}{bar}
\node[op] at (1.48,2.92) {$-$};
\prof{1.75}{2.62}{10,8,6,11,4,7,3,6}{bar}
\node[op] at (3.13,2.92) {$=$};
\prof{3.40}{2.62}{7,4,3,3,3,3,2,2}{bar}

\node[strong] at (0.52,1.94) {shift};
\node[lbl]    at (2.17,1.94) {generic$(c)$};
\node[strong] at (3.95,1.94) {residual $\resid(d,c)$};
\prof{0.10}{1.22}{7,4,3,3,3,3,2,2}{bar}
\node[op] at (1.48,1.52) {$-$};
\prof{1.75}{1.22}{6,4,3,2,3,2,2,2}{faintbar}
\node[op] at (3.13,1.52) {$=$};
\prof{3.40}{1.22}{1,0,0,1,0,1,0,0}{bar}

\node[lbl, anchor=north, align=center] at (2.30,0.82)
  {generic$(c)$ is the mean over the \emph{other}\\drugs in this cell line};

\draw[black!25] (5.05,0.10) -- (5.05,3.60);

% ================================================ RIGHT: the encoding fork
\node[strong, anchor=west] at (5.35,3.44) {the \emph{same} measurement, written two ways};

\node[lbl, anchor=west] at (5.35,2.96) {\textbf{(a)} genes ranked by \emph{expression}};
\node[anchor=west] at (5.35,2.52) {%
  \tikz[baseline]{\node[shared]{ACTB};}\;\tikz[baseline]{\node[shared]{GAPDH};}\;%
  \tikz[baseline]{\node[differs]{CDKN1A};}\;\tikz[baseline]{\node[shared]{TPT1};}\;$\cdots$};
\node[anchor=west] at (5.35,2.06) {%
  \tikz[baseline]{\node[shared]{ACTB};}\;\tikz[baseline]{\node[shared]{GAPDH};}\;%
  \tikz[baseline]{\node[differs]{HMOX1};}\;\tikz[baseline]{\node[shared]{TPT1};}\;$\cdots$};
\node[strong, anchor=west] at (5.35,1.58)
  {only \textbf{34.6 / 200} tokens differ \;(\textbf{17\%})};

\node[lbl, anchor=west] at (5.35,1.00) {\textbf{(b)} genes ranked by \emph{signed residual}};
\node[anchor=west] at (5.35,0.56) {%
  \tikz[baseline]{\node[differs]{CDKN1A};}\;\tikz[baseline]{\node[differs]{MDM2};}\;%
  \tikz[baseline]{\node[gene, fill=black!3]{[DOWN]};}\;\tikz[baseline]{\node[differs]{CCNB1};}\;%
  $\cdots$};
\node[strong, anchor=west] at (5.35,0.08)
  {\textbf{118.2 / 200} differ \;(\textbf{59\%})};

% ================================================ the invariant
\draw[decorate, decoration={brace, amplitude=4pt}, black!55] (0.10,-0.55) -- (11.20,-0.55);
\node[strong, anchor=north] at (5.65,-0.73)
  {and the replicate-retrieval ceiling is \textbf{unchanged}: $0.742 \rightarrow 0.743$};

\end{tikzpicture}
\caption[The residual frame and the encoding fork]{\textbf{The information was always present; the
tokenisation discarded it.} Left: the drug-specific residual is built by two subtractions. The
control removes the cell's own state, leaving the \emph{shift}; the mean over the other drugs in the
same cell line removes the \emph{generic} program that every compound triggers, leaving what makes
this drug different. Right: that same measurement, written two ways. Ranking genes by expression (a)
puts housekeeping genes first, so only $34.6$ of the top $200$ tokens differ between two drugs in the
same context and cross-entropy has almost nothing to learn drug identity from; ranking by signed
residual (b) raises that to $118.2$. The bracket is the point of the figure --- the
replicate-retrieval ceiling is essentially identical under both encodings, so no information has been
added. What changed is how much of it the token sequence exposes to the loss.}
\label{fig:frame}
\end{figure}


\paragraph{Scoring.} Predicted and true residuals are both mapped to a signed rank vector (top-$k$
up positive, top-$k$ down negative, weight $1/\log_2(\textrm{rank}+1)$) so a generated sentence and
a continuous truth are compared like for like; similarity is cosine and the score is $\NIR$ over
other drugs in the same cell line.

\subsection{Generalisation design}
\label{sec:methods-holdout}

Following DrEval's insistence that the split match the intended
application~\cite{dreval}, and using its nomenclature, we build a three-way split keyed at the
\textbf{(drug, cell line)} level so that every plate and dose of a held-out pair goes to the held-out
split and no sibling condition leaks into training:
\begin{description}
  \item[\texttt{unseen\_drug} --- Leave-Drugs-Out (LDO)] every condition of a held-out drug. The
        model never sees the token. The application is drug design. This is the \emph{control}:
        given that chemical structure does not predict response (Section~\ref{sec:res-scope}), it
        should fail, and its failing is what proves that any transfer observed in
        \texttt{unseen\_combo} comes from having seen the drug.
  \item[\texttt{unseen\_combo} --- Leave-Pairs-Out (LPO)] a held-out (drug, cell line) pair whose
        drug is seen in at least three other cell lines. The application is missing-value
        imputation. This is cross-context transfer and the scientifically meaningful test.
  \item[\texttt{train}] everything else.
\end{description}

We do not build a Leave-Cell-Lines-Out (LCO) split here; Section~\ref{sec:limitations} argues it is
the natural next axis and explains why its interest is contingent on the transfer coefficient.

\paragraph{A warning we inherit with the vocabulary.} DrEval reports, citing~\cite{partin}, that
apparently good LPO and LCO performance can be reached simply by predicting each drug's average
response across the training cell lines. Our transfer result is an LPO result and therefore sits
exactly in that regime. This is the reason the baseline ladder below is not optional: the
drug-average predictor is not a formality to be dismissed, it is the hypothesis under test.

\paragraph{Stratified sampling.} Evaluation samples conditions under an explicit per-split quota. A
uniform draw reproduces the pool's proportions, which starves precisely the split the experiment
exists to test: in a first run, $250$ constructed held-out combinations yielded only $33$ scored
conditions, purely as an artifact of the sampler.

\subsection{Baselines}
\label{sec:methods-baselines}

Each is a predictor of the drug-specific residual, scored through an identical pipeline --- same
truths, same signed-rank encoding, same comparison set.

\begin{table}[h]
\centering
\begin{tabular}{lll}
\toprule
Baseline & Information used & Role \\
\midrule
\texttt{drug\_lookup}   & this drug's mean residual, all other cell lines & the bar to beat \\
\texttt{drug\_lookup\_1} & this drug in \emph{one} other cell line & removes the averaging advantage \\
\texttt{moa\_lookup}    & other drugs sharing this drug's MoA, same cell line & mechanism-leak diagnostic \\
\texttt{control\_copy}  & none (predict the control cell) & must be $\approx 0.50$ \\
\texttt{generic}        & none (predict the average response) & must be $\approx 0.50$ \\
\texttt{random}         & none & floor \\
\bottomrule
\end{tabular}
\caption{The baseline ladder. \texttt{drug\_lookup} is the pre-registered win condition: a lookup
can reproduce a memorised average, so only a model that tailors the response to the control cell's
state can beat it.}
\label{tab:baselines}
\end{table}

\texttt{drug\_lookup\_1} selects a cell \emph{line} and averages that line's conditions rather than
selecting a single condition, because $62\%$ of drugs are multi-dose and a single condition would
confound ``different cell line'' with ``different dose''. Baselines are reported per split: on
\texttt{unseen\_drug} the same-drug lookups are \emph{oracles}, since they read the held-out drug's
own residuals, information the model is definitionally denied. \texttt{moa\_lookup} is not an oracle
in any split, as it uses only drugs that were in training.

\subsection{Variance decomposition}
\label{sec:methods-vardecomp}

To separate the per-drug main effect from the drug$\times$cell-line interaction we model
\begin{equation}
  \resid(d,c) \;=\; \beta(d) \;+\; \kappa(d,c) \;+\; \epsilon,
  \qquad
  \Tcoef \;=\; \frac{\operatorname{var}\beta}{\operatorname{var}\beta + \operatorname{var}\kappa},
\end{equation}
where $\beta$ is constant across cell lines and $\kappa$ is the interaction --- the only component a
conditional model can win on.

\paragraph{This is DrEval's naive predictor, in transcriptomic space.} DrEval's
\texttt{NaiveMeanEffectsPredictor}~\cite{dreval} predicts
$\hat{y}_{ij} = \mu + (\mu^c_i - \mu) + (\mu^d_j - \mu) = \mu^c_i + \mu^d_j - \mu$: an additive
model of cell-line and drug main effects, with no interaction term. Our
$\textrm{generic}(c) + \beta(d)$ is precisely that quantity, lifted from a scalar potency to a
$946$-dimensional profile. It follows that $\kappa$ \emph{is} the term their naive model omits, and
that $1 - \Tcoef$ measures exactly the headroom available above it. DrEval finds that deep models
barely beat this additive baseline on drug sensitivity; the transfer coefficient asks whether, in
the transcriptomic setting, there is anything above it to win in the first place.

$\Tcoef$ is estimated by disattenuated cross-line correlation,
\begin{equation}
  \hat{\Tcoef} \;=\; \frac{\cos\!\big(\resid(d,c_1),\, \resid(d,c_2)\big)}
                           {\sqrt{\rho(d,c_1)\,\rho(d,c_2)}},
  \qquad
  \rho \;=\; \frac{2\cos(\resid_A, \resid_B)}{1 + \cos(\resid_A, \resid_B)},
\end{equation}
with $\rho$ the Spearman--Brown reliability of the full-sample residual. Dividing out the
reliability is essential: measurement noise alone depresses the raw cross-line cosine and would
otherwise be read as interaction.

Five controls accompany the estimate, each capable of changing the conclusion:
\begin{enumerate}
  \item \textbf{Control half-split.} The production target builder subtracts one control pseudobulk
        from both halves, so control-estimation error is common to $\resid_A$ and $\resid_B$; this
        inflates the reliability and biases $\Tcoef$ \emph{downward}, manufacturing apparent
        interaction. Both builds are run and the bias reported.
  \item \textbf{Negative control, structure-matched.} A different drug in a \emph{different} cell line, breaking only the drug match while preserving the cross-line structure of the estimand. Must be $\approx 0$. A same-plate/different-drug control was tried first and is \emph{not} adequate: it differs from the estimand in two ways at once, and plate structure surviving a cell-line-scoped generic inflates it to $+0.478$ while leaving the estimand untouched --- which produced a false-alarm rejection before the matched control was built.
  \item \textbf{Simulated null.} At $\kappa = 0$ the estimator does not return $1.0$. A simulation
        calibrated to the measured noise level fixes the reference. In validation the raw cross-line
        cosine at $\kappa = 0$ is $0.734$, which read naively reports $27\%$ interaction where the
        truth is zero.
  \item \textbf{Dose strata.} Same drug, same cell line, different dose is not cross-line transfer;
        it is reported separately and serves as a reference scale.
  \item \textbf{Leave-one-out generic.} The generic includes drug $d$, leaving a $-1/m$ self-term.
\end{enumerate}
Validation against planted ground truth recovers $\Tcoef = 1.00$, $0.70$ and $0.40$ to within
$0.005$, and the different-drug negative control returns $-0.001$.

Finally, $\lVert \resid \rVert / \lVert \shift \rVert$ reports what fraction of the total
perturbation response the drug-specific frame covers --- the number that scopes every claim in this
thesis, and one that has not previously been computed for this task.
